\scp{Labial stops}{12-03}
\citet[133]{Collins1986} observed “the treatment of the labial
stops distinguishes \ili{Banda} from the descendants of Seran Laut”
with which he classified them under the label of “Proto-\ili{Banda}”
(his Seran Laut ≈ our \tsc{Eastern Islands} node of \tsc{Seram-Tanimbar-Bomberai}).
Reflexes of *p for \ili{Banda} are shown in example \qf{12-08}.
This shows the sound correspondences for \ili{Banda} of *p- > {\0},
*-p- > {\0}, *-p = p (one instance).

\noindent{\wg\st{6pt}\tblr{>{\hspace{-\tabcolsep}}l<{\hspace{-\tabcolsep}}
>{\hspace{-\tabcolsep}\enspace}
llll}
\lsptoprule
%\tbx{12-8}	&	PMP	&	\ili{Banda}	&	putative analysis	&	\ili{Banda} gloss	\\	\midrule
%\endfirsthead
\tbx{12-08}		&	PMP	&	\ili{Banda}	&	analysis	&	\ili{Banda} gloss	\\	\midrule
%\endhead
&	*\tbr{p}aRih	&	ari	&		&	stingray	\\
	&	*\tbr{p}ija	&	ilo	&		&	how many, how much	\\
	&	*ha\tbr{p}uy	&	au	&		&	fire	\\
	&	*qa\tbr{p}əju	&	elu-n	&		&	gall, bile	\\
	&	*ni\tbr{p}ay	&	nia, ɲia	&		&	snake	\\
	&	*ma-ni\tbr{p}is	&	munít	&	/muniit/	&	thin 	\\
	&	*ma-la\tbr{p}aR	&	mlar, malar	&	/mlaar/	&	hungry (inflected) 	\\
	&	*sa\tbr{p}an	&	ansá	&	/ansaa/	&	what 	\\
	&	*ə\tbr{p}at	&	at	&	/aat/	&	four 	\\
	&	*qatə\tbr{p}	&	ato\tbr{p}	&		&	thatch, roof	\\
\lspbottomrule
\end{tabular}\mbox{}\\}

Reflexes of *b for \ili{Banda} are shown in example \qf{12-09},
which attests *b > \it{f}.
In the main text of \citet[131--134]{Collins1986} the reflex
of *b before \it{u} is transcribed as a bilabial fricative ‹ɸ›,
while in the word list in the appendix it is transcribed ‹f›.
This suggests allophony in which the labio-dental fricative is
realised as bilabial before /u/; that is /f/ → [ɸ] /{\gap}u.

\noindent{\wg\st{6pt}\tblr{>{\hspace{-\tabcolsep}}l<{\hspace{-\tabcolsep}}
>{\hspace{-\tabcolsep}\enspace}
lll}
\lsptoprule
%\tbx{12-9}	&	PMP	&	\ili{Banda}	&	\ili{Banda} gloss	\\	\midrule
%\endfirsthead
\tbx{12-09}	&	PMP	&	\ili{Banda}	&	\ili{Banda} gloss	\\	\midrule
%\endhead
	&	*\tbr{b}atu	&	\tbr{f}atu	&	rock, stone	\\
	&	*\tbr{b}əRsay	&	\tbr{f}ese	&	paddle	\\
	&	*\tbr{b}aqəRu	&	fer{\tl}\tbr{f}eru-ŋo	&	new, young (also \it{feferúŋ})	\\
	&	*\tbr{b}ulan	&	\tbr{f}ulan	&	moon	\\
	&	*\tbr{b}usuR	&	\tbr{f}usur	&	bow	\\
	&	*\tbr{b}əli	&	fel{\tl}\tbr{f}eli-ko	&	sell (literally `cause to buy')	\\
	&	*\tbr{b}uhək	&	\tbr{f}uk-	&	head hair	\\
	&	*\tbr{b}ulu	&	\tbr{f}ulu-n 	&	feather	\\
	&	*\tbr{b}uaq	&	\tbr{f}ua-n	&	fruit	\\
	&	*\tbr{b}ituqən	&	\tbr{f}otuon	&	star	\\
	&	*\tbr{b}u\tbr{b}uŋ	&	\tbr{f}a\tbr{f}un	&	gable	\\
	&	*\tbr{b}a\tbr{b}aq	&	\tbr{f}o\tbr{f}án	&	below (vowels unexplained)	\\
	&	*\tbr{b}u\tbr{b}u	&	\tbr{f}u\tbr{f}u	&	trap	\\
	&	*qa\tbr{b}aRa	&	\tbr{f}ara-	&	shoulder	\\
	&	*(ma-)\tbr{b}a-\tbr{b}\<in\>ahi	&	mbei\tbr{f}ino	&	female (also \it{meifino} from **ma-binay)	\\
	&	**nti\tbr{b}u	&	ndi\tbr{f}u	&	fly (v.)	\\
	&	*\tbr{b}a-la\tbr{b}aw	&	\tbr{f}ola\tbr{f}a	&	rat	\\
	&	*la\tbr{b}əR	&	la\tbr{f}ár	&	wide	\\
\lspbottomrule
\end{tabular}\mbox{}\\}

Initial,
medial, and final labial stops for \ili{Banda},
\tsc{Eastern Islands} and \tsc{\ili{Teor}-Kur} languages (all \tsc{Seram-Tanimbar-Bomberai})
are shown in \trf{12-04}.
Data are not available for all languages in all positions.

\begin{table}
	\centering\tcp{Labials in \ili{Banda} and nodes of \tsc{Seram-Tanimbar-Bomberai}}{12-04}
	\st{6pt}\begin{tabular}{lllllll}\lsptoprule
PMP				&	*p-				&	*-p-				&	*-p				&	*b-				&	*-b-			&	*bu\\	\midrule
\ili{Bati}			&	{\hr}f		&	\hp{*-}f		&						&	{\hr}b, w	&						&	{\hr}u,wu\\
\ili{Geser}			&	{\hr}f		&	\hp{*-}f		&						&	{\hr}b, w	&	\hp{*-}w	&	{\hr}u\\
\ili{Gorom}			&	{\hr}h		&	\hp{*-}h		&						&	{\hr}w		&	\hp{*-}w	&	{\hr}u,wu\\
\ili{Watubela}	&	{\hr}f		&	\hp{*-}f		&	\hp{*-}f	&	{\hr}w		&						&	{\hr}u\\
\ili{Kesui}			&	{\hr}f		&	\hp{*-}f		&						&	{\hr}b, w	&						&	{\hr}wu\\
\ili{Kowiai}		&	{\hr}ɸ		&	\hp{*-}ɸ		&						&	{\hr}b, ɸ	&	\hp{*-}ɸ	&	{\hr}ɸu\\
\ili{Teor}			&	{\hr}f		&	\hp{*-}f		&						&	{\hr}f		&	\hp{*-}f	&	‹phu›\\
\ili{Kur}				&	{\hr}p		&	\hp{*-}p		&	\hp{*-}p	&	{\hr}f		&	\hp{*-}f	&	{\hr}fu\\	\midrule
\ili{Banda}			&	{\hr}{\0}	&	\hp{*-}{\0}	&	\hp{*-}p	&	{\hr}f		&	\hp{*-}f	&	{\hr}fu [ɸu]\\
		\lspbottomrule
	\end{tabular}
\end{table}

The behaviour of PMP *b in \ili{Banda} is not contrastive compared
with \tsc{Eastern Islands} and \tsc{\ili{Teor}-Kur}.
Collins made his earlier observation based on comparing only
\ili{Banda}, \ili{Geser}, and \ili{Watubela}, for which the reflexes of PMP *b are noticeably different.

While \tsc{Eastern Islands} and \tsc{\ili{Teor}-Kur} languages consistently
retain a reflex of PMP *p in initial
and medial positions, \ili{Banda} consistently loses it.
This could be seen as a secondary development,
but it does set \ili{Banda} apart from those two nodes.
The development of *p in \ili{Banda} is most closely paralleled by
\tsc{Timor-Babar} languages (\crf{ch:TB}), to the south of the \ili{Banda} Islands,
which have initial and medial *p > \it{h} or *p > {\0},
while final *p = **p.
However, apart from this, there does not seem to be additional
evidence for placing \ili{Banda} within \tsc{Timor-Babar}.
